%

\section{Proofs for Impossibility Results}\label{sec:impossibility-proofs}

\subsection{Generalized Assumptions}

First, we restate \cref{assn:inscale-covdist,assn:piecewise} in the multivariate case to allow for the most general version of our results.

\begin{appassumption}\label{assn:inscale-covdist-multivariate}
For any $\covval\in\covspace$, $\featuresub\subseteq[\covdim]$, $\inscale\in\PosReals^{\abs{\featuresub}}$, and $\covdist\in\probspace(\covspace)$, we say the present assumption holds if \cref{assn:inscale-covdist} holds for each $j\in\featuresub$ with $\inscale\feat{j}$.
\end{appassumption}

\begin{appassumption}\label{assn:piecewise-multivariate}
For any $\covval\in\covspace$, $\featuresub\subseteq[\covdim]$, and $\inscale\in\PosReals^{\abs{\featuresub}}$,
let $\covnbhd = \prod_{j\in\featuresub} [\covval\feat{j}-\inscale\feat{j}, \covval\feat{j}+\inscale\feat{j}] \times \covval\feat{[\covdim]\setminus\featuresub}$.
%
For any $(\modeldum\feat{j}:[\covval\feat{j}-\inscale\feat{j}, \covval\feat{j}+\inscale\feat{j}]\to\dataspace)_{j\in\featuresub}$ and $\numpiece\in\Nats$, we say that the present assumption holds with size $\numpiece$ if
\*[
    \Big\{\model\in(\covspace\to\dataspace) \setdelim
    \forall\covvaldum\in\covnbhd \quad \model(\covvaldum) = \sum_{j\in\featuresub}\modeldum\feat{j}(\covvaldum\feat{j}), \ \model\restrict{\comp{\covnbhd}} \in \pwiselin{\numpiece}\restrict{\comp{\covnbhd}}, 
    \ \model \text{ is continuous}\Big\} \subseteq \modelspace.
\]
\end{appassumption}

\subsection{Proof of Main Result}

We state and prove the following multivariate generalization.

\begin{theorem}\label{fact:oracle-hypothesis-multivariate}
Fix any $\covval\in\covspace$, $\featuresub\subseteq[\covdim]$, 
%
$\inscale \in\PosReals^{\abs{\featuresub}}$,
$\covdist \in \probspace(\covspace)$, 
and
$(\modeldum\nullind\feat{j},\modeldum\altind\feat{j}:[\covval\feat{j}-\inscale\feat{j}, \covval\feat{j}+\inscale\feat{j}]\to\dataspace)_{j\in\featuresub}$.
Suppose that \cref{assn:inscale-covdist-multivariate} is satisfied and \cref{assn:piecewise-multivariate} is satisfied with $m=2^{\abs{\featuresub}}$.
%
Define $\covnbhd$ as in \cref{assn:piecewise-multivariate}, and let
%
%
%
%
%
%
\*[
    \modelspace\nullind
    &= \Big\{\model\in\modelspace\setdelim \forall\covvaldum\in\covnbhd, \ \model(\covvaldum) = \sum_{j\in\featuresub}\modeldum\nullind\feat{j}(\covvaldum\feat{j}) \Big\} \\
    \modelspace\altind
    &= \Big\{\model\in\modelspace\setdelim \forall\covvaldum\in\covnbhd, \ \model(\covvaldum) = \sum_{j\in\featuresub}\modeldum\altind\feat{j}(\covvaldum\feat{j})\Big\}.
\]
For any \completename{} and \additivename{} $\expdef$ and \oraclehyptestname{} $\oraclehyptest$, 
\*[
    \specsym_{\expdef,\covdist,\covval}(\oraclehyptest) \leq 1 -
    \senssym_{\expdef,\covdist,\covval}(\oraclehyptest).
\]
\end{theorem}

Before we prove this, we need the following technical result, which is the multivariate analogue of \cref{fact:interior-freedom}.
\begin{theorem}\label{fact:interior-freedom-multivariate}
Fix any $\covval\in\covspace$, $\featuresub\subseteq[\covdim]$, 
%
$\inscale \in\PosReals^{\abs{\featuresub}}$,
$\covdist \in \probspace(\covspace)$, 
and
$(\modeldum\feat{j}:[\covval\feat{j}-\inscale\feat{j}, \covval\feat{j}+\inscale\feat{j}]\to\dataspace)_{j\in\featuresub}$.
Suppose that \cref{assn:inscale-covdist-multivariate} is satisfied and \cref{assn:piecewise-multivariate} is satisfied with $m=2^{\abs{\featuresub}}$.
For every $\expval \in \Reals^{\abs{\featuresub}\times\datadim}$, there exists 
$\model\in\modelspace$ such that for every \completename{} and \additivename{} \fielddefname{}, $\expdef(\model,\covdist,\covval)\feat{\featuresub} = \expval$, and if $\abs[0]{\covval\feat{j}-\covvaldum\feat{j}} \leq \inscale\feat{j}$ for all $j\in\featuresub$ then $\model(\covvaldum) = \sum_{j\in\featuresub} \modeldum\feat{j}(\covvaldum\feat{j})$.
\end{theorem}

We are now able to prove our main theorem.

\begin{proof}[Proof of \cref{fact:oracle-hypothesis-multivariate}]
For $j\not\in \featuresub$, let $\modeldum\nullind\feat{j} = \modeldum\altind\feat{j} \equiv 0$.
Fix $\expval = \mathbf{0}$.
Let $\model\nullind$ be the \modelname{} guaranteed to exist from \cref{fact:interior-freedom-multivariate} for $(\modeldum\nullind\feat{j})_{j\in[\covdim]}$, and similarly define $\model\altind$. Since $\model\nullind,\model\altind\in\modelspace$ by \cref{assn:piecewise-multivariate},
\*[
    \specsym_{\expdef,\covdist,\covval}(\oraclehyptest) 
    &= \inf_{\model\in\modelspace\nullind} [1-\oraclehyptest(\expdef(\model,\covdist,\covval))] \\
    &\leq 1-\oraclehyptest(\expdef(\model\nullind,\covdist,\covval)) \\
    &= 1-\oraclehyptest(\mathbf{0}) \\
    &= 1-\oraclehyptest(\expdef(\model\altind,\covdist,\covval)) \\
    &\leq 1-\inf_{\model\in\modelspace\altind} \oraclehyptest(\expdef(\model,\covdist,\covval)) \\
    &= 1 - \senssym_{\expdef,\covdist,\covval}(\oraclehyptest).
\]
\end{proof}

\subsection{Proof of \cref{fact:interior-freedom-multivariate}}

%
The proof follows by explicitly constructing a piecewise \modelname{} $\model$ that satisfies the desired properties and is linear outside of the neighbourhood around $\covval$.
In particular, for each $j\in\featuresub$ and $k\in[\datadim]$, we construct $\modelleft\feat{jk}$ and $\modelright\feat{jk}$, define
\*[
    \model\feat{jk}(\covvaldum\feat{j})
    =
    \begin{cases}
    \modelleft\feat{jk}(\covvaldum\feat{j})
    &\covvaldum\feat{j} \in (\covleft\feat{j}, \covval\feat{j}-\inscale\feat{j}) \\
    \modeldum\feat{j}(\covvaldum\feat{j})\feat{k}
    &\covvaldum\feat{j} \in [\covval\feat{j}-\inscale\feat{j}, \covval\feat{j}+\inscale\feat{j}] \\
    \modelright\feat{jk}(\covvaldum\feat{j})
    &\covvaldum\feat{j} \in (\covval\feat{j}+\inscale\feat{j},\covright\feat{j}) \\
    0
    & \text{otherwise},
    \end{cases}
\]
and finally define
\*[
    \model(\covvaldum)\feat{k}
    = \sum_{j\in\featuresub} \model\feat{jk}(\covvaldum\feat{j}).
\]
By definition, $\model(\covvaldum) = \sum_{j\in\featuresub} \modeldum\feat{j}(\covvaldum\feat{j})$ if $\abs[0]{\covvaldum\feat{j}-\covval\feat{j}}\leq\inscale\feat{j}$ for all $j\in\featuresub$.
Further, since $\expdef$ is \additivename{}, 
\*[
    \expdef(\model,\covval,\covdist)\feat{jk}
    = \expdef(\model\feat{jk}, \covval\feat{j}, \covdist\feat{j}).
\]
For some $\linparamleft\feat{jk},\linparamright\feat{jk}\in\Reals$ to be chosen in the proof, set
\*[
    \modelleft\feat{jk}(\covvaldum\feat{j})
    = 
    \linparamleft\feat{jk} \cdot (\covvaldum\feat{j}-\covval\feat{j}+\inscale\feat{j}) + \modeldum\feat{j}(\covval\feat{j}-\inscale\feat{j})\feat{k}
\]
and
\*[
    \modelright\feat{jk}(\covvaldum\feat{j})
    = 
    \linparamright\feat{jk} \cdot (\covvaldum\feat{j}-\covval\feat{j}-\inscale\feat{j}) + \modeldum\feat{j}(\covval\feat{j}+\inscale\feat{j})\feat{k}.
\]
These functions are piecewise linear on $(\covleft\feat{j}, \covval\feat{j}-\inscale\feat{j})$ and $(\covval\feat{j}+\inscale\feat{j},\covright\feat{j})$ for each $j\in\featuresub$ respectively, and hence $\model$ is piecewise linear on the product of these intervals (which form convex polytopes). Thus, by \cref{assn:piecewise}, $\model\in\modelspace$.
It remains to construct $\linparamleft\feat{jk}$ and $\linparamright\feat{jk}\in\Reals$ so that $\expdef(\model\feat{jk}, \covval\feat{j}, \covdist\feat{j}) = \expval\feat{jk}$ for each $j\in\featuresub$ and $k\in[\datadim]$.


Fix $j\in[\covdim]$ and $k\in[\datadim]$.
Recall that $\mu_j$ is a probability measure, and hence maps intervals of the form $(a,b)$ to a number in $[0,1]$.

For any $\linparamleft\feat{jk},\linparamright\feat{jk}\in\Reals$ that satisfy
\*[
    &\hspace{-1em}\linparamleft\feat{jk} \Bigg[ \EE_{\covobs\feat{j}\sim\covdist\feat{j}}\Big[\covobs\feat{j} \cdot\ind{\covobs\feat{j} \in (\covleft\feat{j}, \covval\feat{j}-\inscale\feat{j})}\Big] - (\covval\feat{j}-\inscale\feat{j}) \cdot \covdist\feat{j}\Big((\covleft\feat{j}, \covval\feat{j}-\inscale\feat{j})\Big) \Bigg] \\
    &\quad+ \linparamright\feat{jk} \Bigg[ \EE_{\covobs\feat{j}\sim\covdist\feat{j}}\Big[\covobs\feat{j} \cdot\ind{\covobs\feat{j} \in (\covval\feat{j}+\inscale\feat{j},\covright\feat{j})}\Big] - (\covval\feat{j}+\inscale\feat{j}) \cdot \covdist\feat{j}\Big((\covval\feat{j}+\inscale\feat{j},\covright\feat{j})\Big) \Bigg] \\
    &= -\modeldum\feat{j}(\covval\feat{j}-\inscale\feat{j})\feat{k} \cdot \covdist\feat{j}\Big((\covleft\feat{j}, \covval\feat{j}-\inscale\feat{j})\Big)
    -\modeldum\feat{j}(\covval\feat{j}+\inscale\feat{j})\feat{k} \cdot \covdist\feat{j}\Big((\covval\feat{j}+\inscale\feat{j},\covright\feat{j})\Big) \\
    &\quad- \EE_{\covobs\feat{j}\sim\covdist\feat{j}}\Big[\modeldum\feat{j}(\covobs\feat{j})\feat{k} \cdot\ind{\covobs\feat{j} \in[\covval\feat{j}-\inscale\feat{j},\covval\feat{j}+\inscale\feat{j}]}\Big]
    + \modeldum\feat{j}(\covval\feat{j})\feat{k} - \expval\feat{jk},
\]
since $\expdef$ is \completename{},
\*[
    \expdef(\model\feat{jk}, \covval\feat{j}, \covdist\feat{j})
    = \modeldum\feat{j}(\covval\feat{j})\feat{k} - \EE_{\covobs\feat{j}\sim\covdist\feat{j}}[\model\feat{jk}(\covobs\feat{j})]
    = \expval\feat{jk}.
\]

For simplicity, we set either $\linparamleft\feat{jk}$ or $\linparamright\feat{jk}$ to zero (at worst, this inflates the Lipschitz parameter by a factor of 2). 
By assumption (i),
\*[
    \max\Big\{\covdist\feat{j}\Big((\covleft\feat{j}, \covval\feat{j}-\inscale\feat{j})\Big), \covdist\feat{j}\Big((\covval\feat{j}+\inscale\feat{j},\covright\feat{j})\Big) \Big\} > 0,
\]
so one of 
\*[
    \EE_{\covobs\feat{j}\sim\covdist\feat{j}}\Big[\covobs\feat{j} \cdot\ind{\covobs\feat{j} \in (\covleft\feat{j}, \covval\feat{j}-\inscale\feat{j})}\Big] - (\covval\feat{j}-\inscale\feat{j}) \cdot \covdist\feat{j}\Big((\covleft\feat{j}, \covval\feat{j}-\inscale\feat{j})\Big)
\]
and
\*[
    \EE_{\covobs\feat{j}\sim\covdist\feat{j}}\Big[\covobs\feat{j} \cdot\ind{\covobs\feat{j} \in (\covval\feat{j}+\inscale\feat{j},\covright\feat{j})}\Big] - (\covval\feat{j}+\inscale\feat{j}) \cdot \covdist\feat{j}\Big((\covval\feat{j}+\inscale\feat{j},\covright\feat{j})\Big)
\]
are non-zero. 
If 
\*[
    &\hspace{-1em}\abs{\EE_{\covobs\feat{j}\sim\covdist\feat{j}}\Big[\covobs\feat{j} \cdot\ind{\covobs\feat{j} \in (\covleft\feat{j}, \covval\feat{j}-\inscale\feat{j})}\Big] - (\covval\feat{j}-\inscale\feat{j}) \cdot \covdist\feat{j}\Big((\covleft\feat{j}, \covval\feat{j}-\inscale\feat{j})\Big)} \\
    &>
    \abs{\EE_{\covobs\feat{j}\sim\covdist\feat{j}}\Big[\covobs\feat{j} \cdot\ind{\covobs\feat{j} \in (\covval\feat{j}+\inscale\feat{j},\covright\feat{j})}\Big] - (\covval\feat{j}+\inscale\feat{j}) \cdot \covdist\feat{j}\Big((\covval\feat{j}+\inscale\feat{j},\covright\feat{j})\Big)},
\]
set $\linparamright\feat{jk}=0$, and otherwise set $\linparamleft\feat{jk}=0$.
\manualendproof



